\section{Problem Statement}
\label{sec:problem_statement}
Despite these advances, important limitations remain.
Centralized approaches still suffer from the familiar risks of concentration, including service disruption
when the monitoring server fails and performance bottlenecks when a small number of servers must supervise
very large populations of clients
\cite{guoDecentralizedMultiVenueRealTime2025}.
% (Keywords: bottlenecks, centralized, decentralized)
% Guo
% 2025
% Banerjee
% 2022 --> can't find citation
Decentralized approaches avoid some of these issues, but they are not straightforward replacements because failure
detection in asynchronous environments is inherently uncertain: a missing heartbeat or probe response can indicate a
crashed node, but it can also reflect transient delay, packet loss, or local slowness
\cite{nikolicSelfHealingDilemmasDistributed2021} \cite{dadgarLifeguardLocalHealth2018}.
%(Keywords: False positive, heartbeat, latency)
% Nikolic
% 2020
% Dadgar
% 2017
This uncertainty has practical consequences. SWIM achieves strong scalability, but real deployments have shown that
slow message processing can produce false positive failure detections. \cite{dadgarLifeguardLocalHealth2018}.
%(Keywords: False positive)
% Dadgar
% 2017
Lifeguard addresses this problem by incorporating local health awareness, reducing false positives by more than 50× in
controlled tests \cite{dadgarLifeguardLocalHealth2018}.
%(Keywords: Lifeguard, False positive, reduce)
% Dadgar
% 2017
Yet this line of work remains focused primarily on improving failure detection itself, rather than providing an
end-to-end mechanism that also integrates topology recovery or service restoration.
\cite{dasSWIMScalableWeaklyconsistent2002} \cite{dadgarLifeguardLocalHealth2018}.
%(Keywords: False positive, failure detect, fix, restore, recover)
% Das
% 2002
% Dadgar
% 2017

Similarly, self-healing studies often emphasize recovery after failure once some form of local detection has already
occurred. In one decentralized topology-maintenance approach, neighboring nodes detect failures via heartbeat checks
and then coordinate locally to recreate and reconnect the missing node based on stored topology knowledge
\cite{rodriguezDecentralisedSelfHealingApproach2021}.
%(Keywords: stored, recover, restore, reconnect, local)
% Rodr'iguez
% 2020
Other recent systems combine layered fault detection with dynamic reconstruction to preserve continuity under link or
node disruption \cite{guoDecentralizedMultiVenueRealTime2025}.
%(Keywords: reconstruct, reconnect, continuity, preserve, ensure)
% Guo
% 2025
These results are promising, but they also indicate that the literature frequently treats detection, topology knowledge
dissemination, and post-failure healing as separable modules rather than as a unified framework.

Accordingly, the research gap is not the absence of decentralized node health-checking techniques, but the limited
availability of unified end-to-end designs that jointly address neighbor awareness, failure detection, membership
update, and recovery orchestration within one coherent architecture \cite{fuSystematicReviewFuture2025}.
% Fu
% 2024
This motivates the design of a decentralized framework that does not stop at detecting failures, but also supports
coordinated healing actions and sustained system operation after disruption.

% Many health checking systems rely on a centralized structure, where a single central entity collects
% information and makes decisions for the entire network.
% Although this approach is simple to manage, it may introduce a bottleneck and a single point of failure
% \cite{kingCentralizedDecentralizedComputing1983} \cite{patsiasTaskAllocationMethods2023} .
% To address this limitation, this project adopts a decentralized approach.
% Instead of depending entirely on one central controller, nodes participate in monitoring one another and
% responding to failures.

%A dashboard may be used as a visual aid to observe the system state, but the health checking and recovery
% logic remain decentralized.

% \section{...}
% \label{sec:second}
% In This Part You Write The Second Part Of Your Report.

% \subsection{Failure Models}
% \label{subsec:failure_models}
% Explicitly define the failure models considered in this report:
% \begin{itemize}
%     \item \textbf{Crash failure:} a node stops executing and no longer sends responses.
%     \item \textbf{Network partition:} a node is alive but temporarily unreachable from one or more monitors.
%     \item \textbf{Message loss:} heartbeat or health-check messages are dropped by the network.
%     \item \textbf{Delay failure:} a node responds too slowly because of load, garbage collection, or
% network congestion.
%     \item \textbf{Recovery:} a failed or unreachable node becomes healthy again and must rejoin the active
% membership view.
% \end{itemize}
% State which failures are fully handled, partially handled, or outside the project scope.
